% Part: first-order-logic
% Chapter: sequent-calculus
% Section: rules-and-proofs

\documentclass[../../../include/open-logic-section]{subfiles}

\begin{document}

\iftag{FOL}
      {\olfileid{fol}{seq}{rul}}
      {\olfileid{pl}{seq}{rul}}

\olsection{قاعدے تے \usetoken{P}{derivation}}

اگلی بحث لئی، منّ لو کہ $\Gamma, \Delta, \Pi, \Lambda$، !!{sentence}s
دے محدود سلسلے ظاہر کردے نیں۔

\begin{defn}[سیکوئنٹ]
\emph{سیکوئنٹ} ایس روپ دی اک عبارت اے
\[
\Gamma \Sequent \Delta
\]
جتھے $\Gamma$ تے $\Delta$، زبان $\Lang L$ دے !!{sentence}s دے محدود
(ممکنہ طور تے خالی) سلسلے نیں۔ $\Gamma$ نوں \emph{مقدم} آکھیا جاندا
اے، جد کہ $\Delta$ نوں \emph{تالی} آکھیا جاندا اے۔
\end{defn}

\begin{explain}
سیکوئنٹ دے پچھے بدیہی خیال ایہ اے: جے مقدم وچ سارے !!{sentence}s
قائم ہون تاں تالی دے !!{sentence}s وچوں گھٹ توں گھٹ اک قائم ہووے۔ یعنی، جے
$\Gamma = \tuple{!A_1, \dots, !A_m}$ تے
$\Delta = \tuple{!B_1, \dots, !B_n}$، تاں $\Gamma \Sequent \Delta$
اُدوں تے صرف اُدوں قائم ہُندا اے جدوں
\[
(!A_1 \land \cdots \land !A_m) \lif (!B_1 \lor \cdots \lor
!B_n)
\]
قائم ہُندا اے۔ دو خاص صورتاں نیں: اک جتھے $\Gamma$ خالی اے تے دوجی
جدوں $\Delta$~خالی اے۔ جدوں $\Gamma$ خالی اے، یعنی $m = 0$، $\quad
\Sequent \Delta$ اُدوں تے صرف اُدوں قائم ہُندا اے جدوں $!B_1 \lor
\dots \lor !B_n$ قائم ہُندا اے۔ جدوں $\Delta$ خالی اے، یعنی $n = 0$،
$\Gamma \Sequent \quad$ اُدوں تے صرف اُدوں قائم ہُندا اے جدوں
$\lnot(!A_1 \land \dots \land !A_m)$ قائم ہُندا اے۔ اسیں آکھدے آں
کہ کوئی سیکوئنٹ اُدوں تے صرف اُدوں معتبر اے جدوں اوہدے مطابق
!!{sentence} معتبر اے۔
\end{explain}

جے $\Gamma$، !!{sentence}s دا اک سلسلہ اے، تاں اسیں $\Gamma, !A$ اوس
نتیجے نوں لکھدے آں جیہڑا $!A$ نوں~$\Gamma$ دے سجے سرے اُتے شامل کرن
نال ملدا اے (تے $!A, \Gamma$ اوس نتیجے نوں جیہڑا $!A$ نوں~$\Gamma$
دے کھبے سرے اُتے شامل کرن نال ملدا اے)۔ جے $\Delta$ وی
!!{sentence}s دا اک سلسلہ اے، تاں $\Gamma, \Delta$ ایس لکھی ترتیب نال
دونیں سلسلیاں نوں جوڑ کے بنیا سلسلہ اے۔

\begin{defn}[ابتدائی سیکوئنٹ]
\emph{ابتدائی سیکوئنٹ} اک ایہو جہا سیکوئنٹ اے
\iftag{prvFalse,prvTrue}{جیہڑا اگلے روپاں وچوں کسے اک دا ہووے:
  \begin{enumerate}
    \item $!A \Sequent !A$
    \tagitem{prvTrue}{$\quad \Sequent \ltrue$}{}
    \tagitem{prvFalse}{$\lfalse \Sequent \quad$}{}
  \end{enumerate}}
{جیہڑا $!A \Sequent !A$ دے روپ دا ہووے}، زبان دے کسے وی
!!{sentence} $!A$ لئی۔
\end{defn}

سیکوئنٹ حساب وچ !!^{derivation}s سیکوئنٹاں دے کجھ خاص درخت نیں، جتھے
سب توں اُتّے والے سیکوئنٹ ابتدائی سیکوئنٹ ہُندے نیں، تے جے کوئی
سیکوئنٹ اک یا دو ہور سیکوئنٹاں دے تھلّے آوے تاں اوہ استنباط دے کسے
قاعدے راہیں درست طور تے اوہناں توں نکلنا چاہیدا اے۔ $\Log{LK}$ دے
قاعدیاں دیاں دو وڈیاں قسماں نیں:
\emph{منطقی} قاعدے تے \emph{ساختی} قاعدے۔ منطقی قاعدیاں دے ناں ہِٹھلے
سیکوئنٹ وچ $!A$ تے/یا $!B$ والے !!{sentence} دے !!{main
  operator} دے ناں اُتے رکھے جاندے نیں۔ ہر اک دے دو روپ ہُندے نیں:
اک اوس سیکوئنٹ دا استنباط کرن لئی جیہدے کھبے پاسے !!{operator} والا
!!{sentence} ہووے، تے اک اوس لئی جیہدے سجے پاسے ایہ !!{sentence} ہووے۔

\end{document}
